\section{Primitive functions and the excluded two-slope configurations}\label{sec:primitive-grading}
The primitive-function method gives a separate numerical reduction of a hypothetical unequal-character comparison. It describes the Frobenius weights and identifies canonical residual lines without assuming such an identification. The central criterion excludes each of the three configurations at the end of this section. The shorter proof in Section~\ref{sec:central-ramification} does not require this reduction.

\begin{theorem}\label{thm:three-family-reduction}
Let $p\ge5$, $n\ge3$, $p\nmid e$, $(p-1)\mid e$, and $e<p^n$. Let two potentially supersingular height-two groups over the same complete field $K$ have isomorphic $p^n$-torsion and unequal special inertia traces. After completed unramified base change, let $L/K$ be their common tame good field of degree $p^2-1$. Suppose both good models have two Newton slopes. Write their Hasse valuations as $h_i=(p-1)z_i$, normalized by $v_L(L^\times)=\Z$, and let $\alpha_i$ be their tangent descent scalars for a fixed generator of $\operatorname{Gal}(L/K)$. After interchanging the models if necessary, one of the following holds:
\begin{align}
 z_1&=(p-1)p^n+b,&z_2&=p^{n-1}b,
 &p+1\le b\le p^2-p-1,\quad p\nmid b;
 \label{eq:pg-family-one}\\
 z_1&=b+(p-1)p^{n-1}\ell,&z_2&=p^nb,
 &1\le b\le p-1,\quad b+1\le\ell\le p+1.
 \label{eq:pg-family-two}
\end{align}
The third possibility is
\begin{equation}\label{eq:pg-family-sign}
 z_1=z_2=\frac{p-1}{2}p^n,\qquad \alpha_2=-\alpha_1,
\end{equation}
with the further necessary congruence
\begin{equation}\label{eq:pg-family-congruence}
 (p-1)p^{n-2}\mid(p+1)e-\frac{p-1}{2}(p^n+1).
\end{equation}
In this third case the residual representations split into distinct characters, and the given isomorphism interchanges the canonical residual lines with the opposite character lines. Every family is still subject to the strict two-slope inequalities $0<z_i<pe$.
\end{theorem}

The first step is the classical description of a finite additive group by primitive functions \citep[\S\S2.1--2.2]{Koch2015}. Rescaling the torsion graph gives exactly $2n$ primitive dimensions. Counting two Frobenius chains gives a congruence, and a one-dimensional zero-weight space identifies the canonical lines. The resulting families are necessary intermediate configurations and are proved impossible in Corollary~\ref{pg-cor:exclusion}.

\subsection{The rescaled graph and its primitive functions}
We first identify the finite-dimensional object on which the two Frobenius chains can be compared.

Assume $p\ge5$, $n\ge3$, $p\nmid e$, $(p-1)\mid e$ and $e<p^n$. Suppose two potentially supersingular height-two groups have isomorphic $p^n$-torsion and unequal special inertia traces. We consider the two-slope case; Lemma~\ref{lem:one-slope-recovery} treats the other polygons. After completed unramified base change, take the common tame good field $L/K$ of degree $m=p^2-1$, with ring $R$, algebraically closed residue field $k$, and $E=v_L(p)=me$. Choose $\tau,\pi,\zeta$ as in \eqref{rc-eq:hasseweight}. Write
\[
 h_i=(p-1)z_i,\qquad s_i=v_p(z_i),\qquad 0<z_i<pe,
 \qquad i=1,2.
\]
The integrality of $z_i$ follows from the deformation weight. Let $\alpha_i$ be the tangent scalar of $\tau$. The two Frobenius orbits of these scalars are disjoint, and determinant comparison gives
\begin{equation}\label{pg-eq:determinants}
 \alpha_1^{p+1}=\alpha_2^{p+1}.
\end{equation}
We may use exact tame eigenparameters: applying $m^{-1}\sum_{j=0}^{m-1}\alpha_i^{-j}\tau^j$ to a parameter gives one with a unit linear coefficient and exact eigenvalue $\alpha_i$. Such changes preserve all positive point valuations.

\begin{lemma}[Primitive presentation]\label{pg-lem:primitives}
Let $B$ be a finite-dimensional commutative Hopf algebra over $k$, generated by primitive elements $X,Y$. The space
\[
 V=\operatorname{span}_k\{X,X^p,X^{p^2},\ldots,Y,Y^p,Y^{p^2},\ldots\}
\]
is its full primitive space, and $\dim_k B=p^{\dim_k V}$.
\end{lemma}
\begin{proof}
This is the field case of the classical additive-group/primitive-module correspondence; see \citet[\S\S2.1--2.2, pp.~3--4]{Koch2015}. We give the algebra proof needed here. Choose a basis $v_1,\ldots,v_d$ of $V$, write $v_i^p=\sum_j a_{ij}v_j$, and form
\[
 C=k[t_1,\ldots,t_d]/(t_i^p-\textstyle\sum_j a_{ij}t_j).
\]
The monomials with each exponent less than $p$ form a basis. Indeed reduction by the displayed relations decreases total degree; the overlap between the $i$th and $j$th relations reduces to the difference of the two products of their linear right sides, which is zero. Thus reductions have unique normal forms and $\dim C=p^d$.

Each $t_i$ is primitive, and the primitive relations give an additive Hopf structure on $C$. Sending $t_i$ to $v_i$ gives a surjective Hopf map $C\to B$. If its kernel were nonzero, choose a nonzero kernel element $P$ of least normal-form degree $q$. Augmentation and linear independence exclude $q=0,1$. Every term of
\[
 \Delta(P)-P\otimes1-1\otimes P
\]
has degree less than $q$ in each tensor factor. No reduction is needed when expanding a normal monomial, since its individual exponents are less than $p$. The degree-less-than-$q$ subspace injects into $B$ by minimality, so its tensor square also injects. The displayed expression must therefore vanish already in $C\otimes C$.

A primitive normal form with individual exponents less than $p$ is linear. For any monomial of degree at least two, split off one full positive variable exponent if two variables occur; if only one occurs, split off exponent one. The resulting coproduct coefficient is nonzero modulo $p$. Different original monomials have different sums of exponent vectors and cannot cancel that term. Thus $P$ would be linear, a contradiction. The map $C\to B$ is an isomorphism, and the same argument identifies its primitive elements with the span of the $t_i$.
\end{proof}

\begin{lemma}[Minimum torsion values and leading characters]\label{pg-lem:scaling}
Put $D=p^{2n-1}$ and choose $r$ with $rD\equiv1\pmod m$. For model $i$, the minimum nonzero $p^n$-torsion coordinate value is $z_i/D$, attained by $(p-1)D$ points. The rescaled graph has an additive special fibre whose primitive space has dimension $2n$. It carries a $\Z/D\Z$ grading in which $X^{p^j},Y^{p^j}$ have respective weights $-z_1p^j,-z_2p^j$. The tame scalars
\begin{equation}\label{pg-eq:lambda}
 \lambda_i=\alpha_i\zeta^{-rz_i}
\end{equation}
belong to $\F_p^\times$, and no Frobenius iterate of either coordinate is zero. Moreover,
\begin{equation}\label{pg-eq:bounds}
 (p-1)p^{n-1}<z_i<pe<p^{n+1},\qquad s_i\le n.
\end{equation}
\end{lemma}
\begin{proof}
At level one the minimum is $b=z_i/p$, attained by $p^2-p$ noncanonical points. Suppose the level-$j$ minimum is $b/p^{2j-2}$. Each minimum target has all $p^2$ divisions on the low branch, of value $b/p^{2j}$. A larger target value $q<p^2b$ has division value $q/p^2$, strictly larger than that minimum. If $q\ge p^2b$, every division value is at least $b$, also larger. Division points of zero are first-level points and satisfy the same strict inequality. These statements follow from Lemma~\ref{rc-lem:newton}; its range is valid because every nonzero torsion value is at most the maximum first-level value. Induction proves the minimum, its multiplicity $(p^2-p)p^{2n-2}=(p-1)D$, and that the remaining $D$ points, including zero, have larger value.

Let $A$ be the finite free graph algebra. Make the Eisenstein extension $R'=R[\varpi]/(\varpi^D-\pi)$ and, in its generic algebra, put
\[
 X=x/\varpi^{z_1},\qquad Y=y/\varpi^{z_2},\qquad B'=R'[X,Y].
\]
All geometric values of $X,Y$ are integral. Their characteristic polynomials therefore have coefficients in the integrally closed DVR $R'$, so they are integral over $R'$. Consequently $B'$ is finite free, with the same generic algebra as $A$ and rank $p^{2n}$. Rescaling the formal group laws makes every nonlinear term in the first coordinate divisible by $\varpi^{z_1(d-1)}$, where $d\ge2$ is its total degree, and similarly for the second coordinate. The resulting series converge in the finite complete algebras and give integral Hopf operations. Thus $B_0=B'/\varpi B'$ is generated by primitive $X,Y$. Lemma~\ref{pg-lem:primitives} gives primitive dimension $2n$.

The diagonalizable group scheme $\mu_D$ acts by $\varpi\mapsto t\varpi$, fixing $x,y$. It preserves $B'$ and its special fibre and gives the stated grading. This uses the character grading of $\mu_D$, which is valid even though $\mu_D$ is nonreduced in characteristic $p$. Extend the tame generator by $\tau(\varpi)=\zeta^r\varpi$. Its scalars on the normalized coordinates are \eqref{pg-eq:lambda}.

To determine these scalars, perform the same scaling on one entire $p^n$-torsion algebra. Its normalized coordinate has $p^{2n}$ distinct generic values and satisfies a monic degree-$p^{2n}$ polynomial with integral coefficients. The resulting monogenic algebra has that rank and is the whole scaled algebra. Its special polynomial $P_i(U)$ is additive: $P_i(U+V)-P_i(U)-P_i(V)$ vanishes in the tensor square of the quotient, and has degree less than $p^{2n}$ in each variable, so linear independence of normal monomials makes it identically zero. An additive polynomial in characteristic $p$ has only $p$-power exponents, by the binomial-coefficient criterion. The minimum-value count makes its zero root have multiplicity exactly $D$. Hence
\[
 P_i(U)=U^{pD}+c_iU^D,\qquad c_i\ne0.
\]
Its distinct roots form an $\F_p$-line. Tame action preserves that line and scales it by $\lambda_i$, proving $\lambda_i\in\F_p^\times$. A minimum-value point gives a nonzero specialization of each coordinate in the graph, so none of its Frobenius powers vanishes.

The proof of Proposition~\ref{fc-thm:valuation} bounds every nonzero torsion coordinate below by $(p-1)/(p^n-1)$. Applying this to the minimum gives
\[
 z_i\ge\frac{(p-1)p^{2n-1}}{p^n-1}>(p-1)p^{n-1}.
\]
Together with $z_i<pe<p^{n+1}$ this proves \eqref{pg-eq:bounds}.
\end{proof}

\subsection{Counting the two Frobenius chains}
We now obtain divisibility restrictions from the dimension of the primitive space.

\begin{proposition}\label{pg-prop:chains}
If $\lambda_1\ne\lambda_2$, then $s_1=s_2=n$. In every case $\max(s_1,s_2)\ge n-1$. If $\lambda_1=\lambda_2$, order $s=s_1\le s_2=t$ and write $z_1=p^s a$, $z_2=p^t b$ with $p\nmid ab$. Then $a\ne b$ and
\begin{equation}\label{pg-eq:chain-congruence}
 (p-1)p^{\max(0,2n-1-s-t)}\mid a-b.
\end{equation}
\end{proposition}
\begin{proof}
For coordinate $i$, the weights of its successive Frobenius powers are nonzero and distinct until exponent $d_i=2n-1-s_i$, when the weight first becomes zero. These nonzero vectors represent $d_i+1=2n-s_i$ different weights. All powers have the same tame character $\lambda_i$, since it belongs to $\F_p^\times$. If the two characters differ, their chains are in different tame eigenspaces, so $2n\ge(2n-s_1)+(2n-s_2)$. The bounds $s_i\le n$ force both to equal $n$.

Suppose the characters agree and put $d=t-s$. A nonzero weight in the second chain at exponent $j$ can equal a first-chain weight only at exponent $j+d$, by comparing their $p$-adic valuations modulo $D$. Let $j_0$ be the first equality
\[
 z_1p^{j_0+d}\equiv z_2p^{j_0}\pmod D.
\]
It occurs at latest when both weights are zero; all later weights then coincide. The union has $(2n-s)+j_0$ elements, so $j_0\le s$. Thus $p^{2n-1-s}$ divides $z_1p^{t-s}-z_2$, giving the $p$-power part of \eqref{pg-eq:chain-congruence}. Equality of \eqref{pg-eq:lambda}, together with \eqref{pg-eq:determinants}, gives $p-1\mid z_1-z_2$, and hence $p-1\mid a-b$.

If $a=b$, then $z_2=p^d z_1$ and $\alpha_2=\lambda_1\zeta^{rp^dz_1}=\alpha_1^{p^d}$, contradicting disjoint special character orbits. If $t\le n-2$, the nonzero difference in \eqref{pg-eq:chain-congruence} has absolute value at least $(p-1)p^{2n-1-s-t}\ge(p-1)p^{n+1-s}$. But both $a,b$ are positive and strictly less than $p^{n+1-s}$ by \eqref{pg-eq:bounds}. This contradiction proves $t\ge n-1$.
\end{proof}

\begin{proposition}[One prime-to-$p$ Hasse value]\label{pg-prop:prime-families}
If one normalized Hasse value $z_0$ is prime to $p$, the other $z_*$ belongs with it to one of the two families \eqref{eq:pg-family-one} and \eqref{eq:pg-family-two}.
\end{proposition}
\begin{proof}
Proposition~\ref{pg-prop:chains} gives equal leading characters and $v_p(z_*)\in\{n-1,n\}$. Corollary~\ref{rc-cor:strictordering} gives $z_*<z_0$.

If $z_*=p^{n-1}b$, congruence~\eqref{pg-eq:chain-congruence} says that $z_0-b$ is a nonzero multiple of $(p-1)p^n$. Here $b<p^2$, whereas $z_0>(p-1)p^{n-1}>p^2$, so the multiple is positive. Since $z_0<p^{n+1}$, it is exactly the first positive multiple. The lower bound for $z_*$ gives $b>p-1$, hence $b\ge p+1$. Also
\[
 b<\frac{(p-1)p^n}{p^{n-1}-1}
 =p(p-1)+\frac{p(p-1)}{p^{n-1}-1}<p(p-1)+1.
\]
Thus $b\le p^2-p-1$, since the endpoint $p(p-1)$ is divisible by $p$. This proves \eqref{eq:pg-family-one}.

If $z_*=p^n b$, then $1\le b\le p-1$, and the congruence gives $z_0=b+(p-1)p^{n-1}\ell$. The bounds on $z_0$ give $1\le\ell\le p+1$. The strict ordering implies $(p-1)\ell>pb-b/p^{n-1}$, hence $(p-1)\ell\ge pb$ by integrality. For the stated range of $b$, this is equivalent to $\ell\ge b+1$, proving \eqref{eq:pg-family-two}.
\end{proof}

\subsection{Zero weight, canonical lines and lower levels}
We identify when the primitive grading forces the canonical directions to match, and record the quotient congruence needed when they do not.

\begin{lemma}\label{pg-lem:zero-weight}
If the weight-zero part of the primitive space in Lemma~\ref{pg-lem:scaling} has dimension one, the given torsion isomorphism identifies the two canonical $p$-lines.
\end{lemma}
\begin{proof}
Frobenius eventually takes every weight to zero. Its stable image therefore has dimension at most one, and has dimension at least one because neither coordinate is Frobenius nilpotent. For a sufficiently large $N$, the semilinear Fitting decomposition is
\[
 V=\ker(F^N)\oplus\operatorname{im}(F^N).
\]
It follows by stabilization of kernels and images and rank-nullity over the perfect field $k$. The primitive presentation in Lemma~\ref{pg-lem:primitives} splits accordingly as a tensor product. The nilpotent factor has only its zero $k$-point. The one-dimensional bijective factor has equation $u^p=cu$ with $c\ne0$ and exactly $p$ roots. Thus the support of the graph is an additive group of order $p$ in $\mathbf G_a^2$, an $\F_p$-line. Each coordinate is nonzero somewhere on that support, so both coordinate projections are injective.

Consequently the two rescaled coordinates vanish on specialization at exactly the same generic graph points. By the minimum-value induction, a coordinate vanishes precisely when the $p^{n-1}$ multiple of that point is on the canonical $p$-line: the minimum points are exactly the successive low divisions of the noncanonical first-level points. The torsion isomorphism commutes with multiplication by $p^{n-1}$, which is surjective onto $p$-torsion. Equality of these vanishing kernels identifies the canonical lines.

In particular, the criterion applies when the union of the two chains already represents $2n$ weights. Every represented weight is nonzero as a vector; the primitive dimension is $2n$, so weight zero then has dimension one.
\end{proof}

\begin{lemma}[A lower-level congruence]\label{pg-lem:lower-level}
For an unequal-type two-slope comparison at level $p^N$, $N\ge1$, suppose the normalized Hasse values are $z$ and $p^t b$, with $p\nmid zb$. No bound on $e$ is needed. Then
\begin{equation}\label{pg-eq:lower-level}
 (p-1)p^{\max(0,2N-1-t)}\mid z-b.
\end{equation}
\end{lemma}
\begin{proof}
Rescale with $D_N=p^{2N-1}$. The prime-to-$p$ coordinate has $2N$ distinct weights, including zero, so its chain spans the whole primitive space. The second coordinate has the same leading tame character, and its weight must appear in that chain. If $t<2N-1$, its only possible nonzero matching weight is at exponent $t$, giving $zp^t\equiv p^tb\pmod{D_N}$. If $t\ge2N-1$, no $p$-power divisibility is required. Equality of the leading and determinant characters always gives $p-1\mid z-p^tb$, hence $p-1\mid z-b$. The argument is valid also at $N=1$: it uses only the primitive dimension and grading, not the bound on $e$ in \eqref{pg-eq:bounds}.
\end{proof}

Let $C_i$ denote the canonical residual line, with character $\kappa_i$, and let $\eta_i$ be the quotient character. The high-root Hensel argument of Lemma~\ref{fc-lem:rational} and determinant comparison make both characters trivial over $G_L$. On the tame generator,
\begin{equation}\label{pg-eq:residual-characters}
 \kappa_i(\tau)=\zeta^{(p+1)e-z_i}\alpha_i^{-1},
 \qquad \eta_i(\tau)=\lambda_i^{-1}=\zeta^{pz_i}\alpha_i^{-1}.
\end{equation}
The first formula is \eqref{rc-eq:kappa}. For the second, a minimum-value $p^n$-point has leading field multiplier $\zeta^{rz_i}$; descended point action divides this by $\alpha_i$. Its leading quotient is the noncanonical residual quotient by the minimum-point description. Since $D\equiv p\pmod m$ and $p^2\equiv1\pmod m$, we have $r\equiv p\pmod m$. This proves the formula, independently of the lift of $\tau$, since the quotient character is trivial over $G_L$. It gives
\begin{equation}\label{pg-eq:scalar-condition}
 \kappa_i=\eta_i\quad\Longleftrightarrow\quad z_i\equiv e\pmod{p-1}.
\end{equation}
If the $\lambda_i$ agree, the quotient characters agree, and determinant equality makes the canonical characters agree. Failure of canonical-line matching then forces these two common characters to coincide: otherwise the canonical-character eigenspace is unique. In the coincident case, the transported noncanonical stable line together with the canonical line makes each residual representation scalar.

We will also use the following consequence of the quotient formulas and markings of Proposition~\ref{rc-prop:quotients} and Lemma~\ref{fc-lem:rational}. In normalized values, the canonical quotient changes $z$ to $pz$ for $z<e$, to $(p+1)e-z$ for $z>e$, and to one slope for $z=e$. Its dual kernel is noncanonical in the first case and canonical in the second. A noncanonical stable quotient changes $z$ to $z/p$ and has canonical dual kernel. A paired quotient of a level-$p^n$ comparison retains level $p^{n-1}$ and the dual marking.

Call $z>e$ the central case and $z<e$ the outer case. If $p\mid z$, equality $z=e$ is impossible. A central canonical quotient has prime-to-$p$ Hasse value and nontrivial residual wild action with a unique canonical fixed line. An outer canonical quotient has a stable noncanonical dual line, so its residual action over $L$ is trivial. Hence matched canonical quotients cannot pair a central and an outer source when the central source has $p$-divisible Hasse value. Two central sources with $p$-divisible values have prime-to-$p$ canonical quotient values, so Theorem~\ref{rc-thm:primehasse} gives equal quotient types. The isogenies are defined over $K$ and preserve rational special inertia traces. We have therefore excluded a matched-canonical comparison with both Hasse values divisible by $p$ whenever at least one source is central.

\subsection{Eliminating positive Hasse valuations}
We now prove that the only configuration with two positive $p$-adic Hasse valuations is the quadratic-sign configuration in the main reduction.

\begin{proposition}\label{pg-prop:positive}
If $s_1,s_2>0$, then \eqref{eq:pg-family-sign} and \eqref{eq:pg-family-congruence} hold. The residual representations are split with distinct characters, and the torsion isomorphism interchanges their canonical and opposite residual lines.
\end{proposition}
\begin{proof}
Order $s=s_1\le s_2=t$. Proposition~\ref{pg-prop:chains} leaves $t=n-1$ or $n$.

First suppose $t=n-1$ and $s<n-1$. The leading characters agree. Write $z_1=p^sa$, $z_2=p^{n-1}b$. Congruence~\eqref{pg-eq:chain-congruence} gives
\[
 z_1=p^sb+(p-1)p^n\ell.
\]
Here $\ell\ne0$. A negative value would force $b-a\ge(p-1)p^{n-s}$, although $b<p^2$ and $n-s\ge2$. Thus $\ell>0$, and $z_1<p^{n+1}$ makes $\ell=1$. The lower bound on $z_2$ gives $b\ge p+1$. Both $z_i>p^n>e$ are central. The first chain coincidence is at $j_0=s$, because $v_p(a-b)=n-s$; thus the weight union has $2n$ elements. Lemma~\ref{pg-lem:zero-weight} matches the canonical lines, and the central-quotient exclusion above is a contradiction.

If $s=t=n-1$, the nonzero difference $a-b$ is divisible by $(p-1)p$. But the bounds and prime-to-$p$ condition give $p+1\le a,b\le p^2-1$, so $|a-b|<(p-1)p$. This is impossible. All cases with $t=n-1$ are excluded.

Now suppose $t=n$ and $1\le s<n$. Write $z_2=p^nb$, where $1\le b\le p-1$. The leading characters again agree, and
\begin{equation}\label{pg-eq:mixed-positive}
 z_1=p^sb+(p-1)p^{n-1}\ell,\qquad1\le\ell\le p+1.
\end{equation}
The lower bound on $z_1$ gives positivity, and its upper bound gives the displayed range. If $p\nmid\ell$, then $j_0=s$, so the canonical lines match. Model two is central; the matched-canonical exclusion applies whether model one is central or outer. Thus $\ell=p$, the only multiple of $p$ in its range. Both sources are now central. Their canonical lines cannot match, so equality of the leading characters forces scalar residual action by \eqref{pg-eq:scalar-condition}. Since $(p-1)\mid e$ and $z_i\equiv b\pmod{p-1}$, we obtain $b=p-1$. Hence
\[
 z_1=(p-1)(p^n+p^s),\qquad z_2=(p-1)p^n.
\]
Pair the canonical line of model two with its noncanonical inverse image in model one. The quotient values, at remaining level $p^{n-1}$, are
\[
 w_1=(p-1)(p^{n-1}+p^{s-1}),\qquad
 w_2=(p+1)e-(p-1)p^n.
\]
The second is prime to $p$. If $s=1$, both are, contradicting Theorem~\ref{rc-thm:primehasse}. If $s\ge2$, write $e=(p-1)u$ and apply Lemma~\ref{pg-lem:lower-level} at $N=n-1$; it gives
\begin{equation}\label{pg-eq:positive-arithmetic}
 p^{2n-s-2}\mid A_0:=(p+1)u-p^n-1-p^{n-s}.
\end{equation}
The two-slope and base bounds give
\[
 p^{n-1}+p^{s-1}<u<\frac{p^n}{p-1},\qquad
 0<A_0<\frac{2p^n}{p-1}<3p^{n-1}.
\]
For $s\le n-2$, the divisor in \eqref{pg-eq:positive-arithmetic} is at least $p^n$, whereas $A_0<2p^n/(p-1)<p^n$. For $s=n-1$, integrality gives $u\ge(p+1)p^{n-2}+1$, and therefore $A_0\ge(2p+1)p^{n-2}>2p^{n-1}$. There is no multiple of $p^{n-1}$ strictly between $2p^{n-1}$ and $3p^{n-1}$. Both cases contradict \eqref{pg-eq:positive-arithmetic}.

It remains that $s=t=n$. Write $z_i=b_ip^n$, $1\le b_i\le p-1$. Both sources are central, so their canonical lines cannot match. Pairing each canonical line with its noncanonical image in the other model and applying Lemma~\ref{pg-lem:lower-level} at level $n-1$ gives
\[
 \begin{split}
 (p-1)p^{n-2}&\mid(p+1)e-b_1p^n-b_2,\\
 (p-1)p^{n-2}&\mid(p+1)e-b_2p^n-b_1.
 \end{split}
\]
Subtracting gives $p^{n-2}\mid b_1-b_2$. Since $n\ge3$ and $|b_1-b_2|<p$, we have $b_1=b_2=b$. Reduction modulo $p-1$ gives $(p-1)\mid2b$, so $b=(p-1)/2$ or $p-1$. The latter would make both residual representations scalar. Their scalar characters, and hence their $\lambda_i$, would agree; equal $z_i$ in \eqref{pg-eq:lambda} would then imply $\alpha_1=\alpha_2$, contrary to unequal type.

Thus $b=(p-1)/2$. Write $\alpha_i=\zeta^{a_i}$. The deformation weight gives $z_i\equiv-a_i\pmod{p+1}$. Equal $z_i$ makes $\alpha_2/\alpha_1\in\mu_{p-1}$; determinant equality also puts it in $\mu_{p+1}$. Its order divides two. Equality is excluded by unequal type, so $\alpha_2=-\alpha_1$. The first quotient congruence is exactly \eqref{eq:pg-family-congruence}. Equation~\eqref{pg-eq:scalar-condition} now gives distinct residual characters. A transported noncanonical stable line splits the representation, and the failure of canonical matching forces interchange with the opposite character line. This proves every assertion.
\end{proof}

\begin{proof}[Proof of Theorem~\ref{thm:three-family-reduction}]
If a normalized Hasse value is prime to $p$, Proposition~\ref{pg-prop:prime-families} gives the first two families. If both are divisible by $p$, Proposition~\ref{pg-prop:positive} gives the third. The graph grading and quotient argument provide necessary conditions only. They neither construct a comparison in any family nor identify groups with opposite tangent characters as quadratic twists.
\end{proof}

\begin{corollary}\label{pg-cor:exclusion}
None of the three configurations in Theorem~\ref{thm:three-family-reduction} can occur for an unequal-character comparison under its hypotheses.
\end{corollary}
\begin{proof}
In the first family, $z_2=p^{n-1}b\ge(p+1)p^{n-1}>p^n>e$. In the second, $z_2=p^nb\ge p^n>e$. In the third, $z_1=z_2=((p-1)/2)p^n\ge2p^n>e$. Thus every family has a central model with $p$-divisible normalized Hasse value $z$. Since $p\nmid e$, the integer $(p+1)e-z$ is prime to $p$. The strict two-slope bound $z<pe$ is part of the hypotheses, and the available torsion level is at least $p^2$. Theorem~\ref{thm:central-p2-rigidity} therefore forces equal special traces, a contradiction in every case.
\end{proof}

For instance, the arithmetic quadratic-sign candidates at $p=5,n=3$ have $z_1=z_2=250$ and $e=52,72,92,112$. Each is central and has $5\nmid6e-250$, so the same criterion excludes all four. This exclusion uses actual point fields; it makes no identification of arbitrary opposite-sign groups as quadratic twists.
